\documentclass[11pt]{amsart}

\usepackage[T1]{fontenc}
\usepackage{lmodern}
\usepackage{microtype}
\usepackage{amsmath,amssymb,amsthm}
\usepackage[hidelinks]{hyperref}

\newtheorem{theorem}{Theorem}

\title[A counterexample to Mili\'cevi\'c's conjecture]{A Counterexample at \texorpdfstring{$n=6$}{n=6} to Mili\'cevi\'c's Blocking-Point Conjecture}
\date{}
\subjclass[2020]{Primary 52C10; Secondary 52C30, 51A20, 14H50}
\keywords{blocking set, points in general position, spanned line,
cubic curve, plane configuration, counterexample}

\hypersetup{
  pdftitle={A Counterexample at n=6 to Milicevic's Blocking-Point Conjecture}
}

\begin{document}

\begin{abstract}
Mili\'cevi\'c conjectured that if $n$ points in general position are
blocked by $n$ additional points, then all $2n$ points lie on a cubic
curve. We give an explicit configuration defined over
$\mathbb{Q}(\sqrt5)$ with $n=6$ that satisfies the blocking hypothesis
but lies on no cubic. Together with the theorem of Bhore and Swanepoel
for $n=5$, this shows that $n=6$ is the least possible counterexample.
\end{abstract}

\maketitle

\section{The counterexample}

Let $B$ be a set of $n$ points in the real affine plane, no three
collinear, and let $R$ be a disjoint set of $n$ points such that every
line spanned by two points of $B$ contains a point of $R$.
Mili\'cevi\'c conjectured that $B\cup R$ lies on a cubic curve
\cite[Conjecture~5]{Milicevic}. Keller and Pinchasi proved the conjecture
under the additional requirement that every blue pair has a blocker
outside the segment joining it
\cite[Conjecture~1.2 and Theorem~1.3]{KellerPinchasi}. Bhore and
Swanepoel proved the case $n=5$ \cite[Theorem~2.9]{BhoreSwanepoel}.
For $n\leq4$, the conclusion follows from the ten-dimensionality of the
space of homogeneous plane cubics. Thus the following theorem treats
the first possible value of $n$. Throughout, a cubic is the zero set in
$\mathbb{P}^2$ of a nonzero homogeneous polynomial of degree three;
reducible and nonreduced cubics are included.

\begin{theorem}\label{thm:main}
There exist disjoint six-point sets $B$ and $R$ in the real affine
plane such that no three points of $B$ are collinear and every line
spanned by two points of $B$ contains a point of $R$, but $B\cup R$ is
contained in no cubic curve.
\end{theorem}

\begin{proof}
Set
\[
 p=\frac{3-\sqrt5}{2}.
\]
Then $p^2=3p-1$ and $0<p<\tfrac12$. In $\mathbb{P}^2(\mathbb{R})$, let
\[
\begin{aligned}
 b_1&=(1:0:0), & b_2&=(0:1:0), & b_3&=(0:0:1),\\
 b_4&=(1:1:1), & b_5&=(p:1-p:1), & b_6&=(1-p:p:1),
\end{aligned}
\]
and
\[
\begin{aligned}
 r_1&=(p:0:1), & r_2&=(1-p:1:1), & r_3&=(1:1-p:1),\\
 r_4&=(0:p:1), & r_5&=(-1:-1:0), & r_6&=(1:1:2).
\end{aligned}
\]
The three points with $z=0$ are distinct. After normalization to
$z=1$, the remaining nine points are also distinct because
$0<p<\tfrac12$. Hence all twelve points are distinct. Moreover,
$x+y+z$ takes only the nonzero values
\[
 1,\ 2,\ 3,\ 1+p,\ 3-p,\ -2,\ 4
\]
on these points. We therefore regard
$\mathbb{P}^2\setminus\{x+y+z=0\}$ as a common affine chart.

For $i<j<k$, write
\[
 [ijk]=\det(b_i,b_j,b_k),
\]
with the coordinate vectors as columns. Reduction by $p^2=3p-1$
gives the complete list
\[
\begin{array}{c|l}
\text{value of }[ijk] & \text{triples }ijk\\ \hline
 1 & 123,124,125,126,234\\
 -1 & 134\\
 p & 145,235\\
 -p & 136,246\\
 1-p & 146,236\\
 -(1-p) & 135,245\\
 1-2p & 156,256,345,456\\
 -(1-2p) & 346,356
\end{array}
\]
All these values are nonzero, so no three points of $B$ are collinear.

Write $ij$ for the line through $b_i$ and $b_j$. Direct evaluation of
$\det(b_i,b_j,r_k)$ gives exactly the following incidences:
\[
\begin{array}{c|l}
\text{red point} & \text{blue pairs blocked}\\ \hline
 r_1 & 13,25,46\\
 r_2 & 14,26,35\\
 r_3 & 15,24,36\\
 r_4 & 16,23,45\\
 r_5 & 12,34\\
 r_6 & 34,56
\end{array}
\]
Every one of the fifteen blue pairs occurs, so $R$ blocks all lines
spanned by $B$. The stronger hypothesis of Keller and Pinchasi is not
satisfied: using the normalized representatives with $x+y+z=1$, one has
\[
 r_5=\frac{b_1+b_2}{2},
 \qquad
 r_6=\frac{b_5+b_6}{2},
\]
and the incidence table shows that these are the only red points on
$12$ and $56$, respectively.

It remains to exclude a cubic. Suppose that a nonzero homogeneous
cubic $F(x,y,z)$ vanishes at all twelve points. The four distinct points
\[
 b_3,\quad b_4,\quad r_5,\quad r_6
\]
lie on the line $x-y=0$. The restriction of $F$ to this line is a
binary cubic with four distinct zeros, hence is identically zero.
Therefore
\[
 F=(x-y)Q,
 \qquad
 Q=Ax^2+By^2+Cz^2+Dxy+Exz+Gyz.
\]
The points $b_1,b_2,r_1,r_4,b_5,r_2$ lie off $x-y=0$, so $Q$ vanishes
at each of them. Evaluation at $b_1$ and $b_2$ gives $A=B=0$, while
evaluation at $r_1$ and $r_4$ gives
\[
 C+pE=0,
 \qquad
 C+pG=0.
\]
If $C=0$, then $E=G=0$, and evaluation at $b_5$ gives
$Dp(1-p)=0$, hence $D=0$, contradicting $F\neq0$. Thus $C\neq0$.
After scaling, take $C=p^2$. Then $E=G=-p$, and evaluation at $b_5$
gives
\[
 0=Q(b_5)=p(1-p)(D-1),
\]
so $D=1$. Consequently,
\[
 Q=xy-pxz-pyz+p^2z^2=(x-pz)(y-pz).
\]
At $r_2=(1-p:1:1)$, however,
\[
 F(r_2)=(-p)(1-2p)(1-p)\neq0,
\]
a contradiction. Hence no cubic contains $B\cup R$.
\end{proof}

\begin{thebibliography}{9}

\bibitem{BhoreSwanepoel}
S.~Bhore and K.~Swanepoel,
\emph{On sets of monochromatic objects in bicolored point sets},
\href{https://arxiv.org/abs/2602.17637}{arXiv:2602.17637}, 2026.

\bibitem{KellerPinchasi}
C.~Keller and R.~Pinchasi,
\emph{On sets of $n$ points in general position that determine lines
that can be pierced by $n$ points},
Discrete Comput. Geom. \textbf{64} (2020), no.~3, 905--915.
\href{https://doi.org/10.1007/s00454-020-00201-3}
{doi:10.1007/s00454-020-00201-3}

\bibitem{Milicevic}
L.~Mili\'cevi\'c,
\emph{Classification theorem for strong triangle blocking arrangements},
Publ. Inst. Math. (Beograd) (N.S.) \textbf{107}(121) (2020), 1--36.
\href{https://doi.org/10.2298/PIM2021001M}
{doi:10.2298/PIM2021001M}

\end{thebibliography}

\end{document}